%-------------------------------------------------------------------------------
% Technical appendices, English translation of the source appendices.
%-------------------------------------------------------------------------------

\section{Knowledge-Graph Relation Types and Retrieval Admissibility}
\label{app:kg}

The Deterministic Context Retrieval procedure of \S\ref{sec:method-graphrag} builds context around a set of
core relations and explicitly distinguishes the relations admissible into
retrieval from the relations used only in post-processing,
as shown in Table~\ref{tab:kgrel}. The retrieval stage
obeys four constraints, GR-1 through GR-4. No inferred normative rule may enter
retrieval; only raw normative-source nodes may enter retrieval; no inferential
relation may enter retrieval; and all context must be traceable to a
normative source.

\begin{table}[!htbp]
\centering
\footnotesize
\setlength{\tabcolsep}{4pt}
\caption{Knowledge-graph relation types and retrieval admissibility.}
\label{tab:kgrel}
\begin{tabularx}{\columnwidth}{@{}l l X c@{}}
\toprule
Type & Relation & Role & In retrieval \\
\midrule
core & CONTAINS & containment hierarchy among normative sources, sections, normative rules & Yes \\
core & DEFINES & link from a term/field/definition passage to its defined
  object & Yes \\
core & REFERENCES & explicit reference to another normative source/section/algorithm & Yes \\
core & APPLIES TO & link from a normative rule to the abstract concept it constrains
  & Yes \\
aux & AFFECTS & link from a normative rule to the concrete certificate field it affects;
  used by the field resolver & No \\
aux & DERIVED FROM & provenance link from a derived normative rule to its source section
  & No \\
inferred & OVERRIDES & a more specific normative rule overrides a general one, produced by
  the rule engine & No \\
inferred & CONFLICTS WITH & potential conflict between normative rules, produced by the
  rule engine & No \\
\bottomrule
\end{tabularx}
\end{table}

\section{Key Fields of the Structured IR}
\label{app:ir}

\begingroup\emergencystretch=2em\relax
The system internally uses a JSON-serialized, machine-readable IR object.
Beyond the core four-tuple of \S\ref{sec:method-dual-layer}, its key fields are listed in Table~\ref{tab:irfields};
the remaining fields support provenance, normalization, matching, and downstream
product generation. The lintability determination in \S\ref{sec:lintability} directly
uses only the four fields obligation, assertion\_subject, enforcement\_phase, and
rule\_category.
\par\endgroup


\begin{table}[!htbp]
\centering
\footnotesize
\caption{Key fields of the structured IR.}
\label{tab:irfields}
\begin{tabularx}{\columnwidth}{@{}l l X@{}}
\toprule
Field & Role & Meaning \\
\midrule
subject & core & constrained certificate field or path \\
obligation & core & RFC~2119 deontic level, such as MUST/SHOULD/\dots \\
predicate & core & constraint relation, such as presence, equality,
  containment, or range \\
constraint & core & constraint value, pattern, or condition \\
rule\_text & key ext. & verbatim text of the normative rule sentence, bound
  unchanged to the Description metadata of the generated lint \\
rule\_category & key ext. & normative rule's semantic category, such as
  encoding\_constraint, definition, algorithm\_ref, \dots \\
assertion\_subject & key ext. & assertion subject, such as Certificate / CA /
  RelyingParty / external ecosystem \\
enforcement\_phase & key ext. & phase on which the constraint depends,
  such as encoding / runtime / external validation \\
source\_section & key ext. & normative source and section
  identifier \\
source\_span & key ext. & character/sentence span in the source text \\
evidence\_text & key ext. & source fragment supporting this IR record \\
context\_nodes & key ext. & definition, field metadata, or cited
  context nodes from Deterministic Context Retrieval \\
\bottomrule
\end{tabularx}
\end{table}

\section{Construction Details of the Restricted Code Space \texorpdfstring{$\Tv$}{T V}}
\label{app:tv}

\textbf{The seven components of the vocabulary $\mathcal{V}$.} $\mathcal{V}$ is
the disjoint union of seven finite sets
shown in Table~\ref{tab:vocab}, namely certificate fields $\mathcal{F}_{\mathrm{cert}}$,
DN fields $\mathcal{F}_{\mathrm{dn}}$ over the Subject and Issuer names, OID
constants $\mathcal{O}$, including zlint util.* extension and algorithm OIDs
together with well-known CABF and RFC OIDs, KeyUsage bits $\mathcal{B}_{\mathrm{KU}}$,
ExtKeyUsage bits $\mathcal{B}_{\mathrm{EKU}}$, ASN.1 encoding types
$\mathcal{E}_{\mathrm{ASN1}}$, and named regexes $\mathcal{R}_{\mathrm{regex}}$,
each an audited literal RE2 pattern. Table~\ref{tab:vocab} gives the size and
representative members of each component.
All components are finite sets; the LLM receives a full enumeration of
each component in the prompt, so any identifier outside $\mathcal{V}$ in its
output triggers a parse error.

\begin{table*}[tb]
\centering
\footnotesize
\caption{The seven frozen components of the restricted vocabulary $\mathcal{V}$,
with their sizes and representative members. The LLM may emit only identifiers
drawn from these closed sets.}
\label{tab:vocab}
\begin{tabularx}{\textwidth}{@{}l l r X@{}}
\toprule
Symbol & Component & Size & Representative members \\
\midrule
$\mathcal{F}_{\mathrm{cert}}$ & Certificate fields & 52 &
  c.Version, c.NotAfter, c.KeyUsage, c.DNSNames, c.IPAddresses \\
$\mathcal{F}_{\mathrm{dn}}$ & DN fields over Subject / Issuer & 38 &
  Subject.CommonName, Subject.Organization, Subject.Country, Issuer.CommonName \\
$\mathcal{O}$ & OID constants & 99 &
  util.BasicConstOID, util.SubjectAlternateNameOID, util.AiaOID \\
$\mathcal{B}_{\mathrm{KU}}$ & KeyUsage bits & 9 &
  DigitalSignature, KeyEncipherment, CertSign, CRLSign \\
$\mathcal{B}_{\mathrm{EKU}}$ & ExtKeyUsage bits & 11 &
  ServerAuth, ClientAuth, CodeSigning, EmailProtection \\
$\mathcal{E}_{\mathrm{ASN1}}$ & ASN.1 encoding types & 11 &
  PrintableString, IA5String, UTF8String, BMPString \\
$\mathcal{R}_{\mathrm{regex}}$ & Named RE2 regexes & 28 &
  Re\_LDH\_Hostname, Re\_Rfc5321Mailbox, Re\_TorV3Onion \\
\bottomrule
\end{tabularx}
\end{table*}

\textbf{Representative atomic templates.} The atomic-template set $\mathcal{A}$
has $|\mathcal{A}|=183$ and is organized by semantic cluster, each with a typed
signature $\mathrm{sig}(a)$ described in \S\ref{subsec:tv}. Table~\ref{tab:atoms} excerpts
key atomic templates by cluster.

\begin{table*}[tb]
\centering
\footnotesize
\caption{Representative atomic templates by semantic cluster.}
\label{tab:atoms}
\begin{tabularx}{\textwidth}{@{}l p{4.2cm} l X@{}}
\toprule
Cluster & Atomic template & Signature & Semantics \\
\midrule
I Extension presence & ExtPresent / ExtCritical &
  $(\mathcal{O})$ & extension OID present / present and critical bit set \\
II Macro attributes & IsCA / KeyUsageHas /
  ExtKeyUsageHas & $()$ / $(\mathcal{B}_{\mathrm{KU}})$ /
  $(\mathcal{B}_{\mathrm{EKU}})$ & BasicConstraints.cA$=$true / KeyUsage bit /
  EKU item \\
III Field value/form & FieldEq / FieldMatchesRegex /
  FieldEncodedAs & $(\mathcal{F}, \cdot)$ & field equals literal /
  matches named regex / ASN.1 encoding type \\
IV List/byte-level & ListAllMatch / OidListContains /
  BytesContainsOidDer & $(\mathcal{F}, \cdot)$ & list all satisfy
  subtree / contains named OID / bytes contain the OID's DER \\
V NameConstraints & SubtreeIPListAny\-HasOctetCount\-AndNotAllZero &
  $(\mathcal{F}, \mathbb{Z})$ & NC IP subtree has an entry of given octet count
  and not all-zero \\
VI Special structure & DomainComponentOrdered & $()$ &
  domainComponent in reverse order per RFC~4519 \\
\bottomrule
\end{tabularx}
\end{table*}

There are three combinators $\{\neg, \wedge, \vee\}$, with semantics following
classical propositional logic.

\textbf{Generality grading of atomic templates.} The GENERIC / SPECIALIZED
grading is defined in \S\ref{subsec:atoms}; of the 183 atomic templates in
$\mathcal{A}$, 118 are GENERIC and 65 are
SPECIALIZED. Table~\ref{tab:specialized} lists only representative SPECIALIZED
examples.

\begin{table*}[tb]
\centering
\footnotesize
\caption{Representative SPECIALIZED atomic templates.}
\label{tab:specialized}
\begin{tabularx}{\textwidth}{@{}p{4.6cm} p{3cm} X@{}}
\toprule
Representative SPECIALIZED atomic template & Specialization target &
  Semantics \\
\midrule
SigAlgMatchesTBSSignature & certificate signature-algorithm fields &
  outer signatureAlgorithm byte-for-byte identical to
  tbsCertificate.signature \\
NotAfterIsNoExpirySentinel & validity sentinel value & whether
  notAfter is a specific no-expiry sentinel time \\
DomainComponentOrdered & Subject DN's domainComponent &
  whether the domainComponent sequence is ordered per RFC~4519 \\
CertPolicy\-ExplicitText\-HasEncodingTagInSet & certificatePolicies
  explicitText & whether explicitText uses an allowed ASN.1 string encoding \\
CRLDPHasNameRelative & CRL Distribution Points & whether
  distributionPoint uses the nameRelativeToCRLIssuer form \\
SubtreeIPListAny\-HasOctetCount\-AndNotAllZero & NameConstraints IP subtree
  & whether the IP subtree has an entry of given octet count that is not
  all-zero \\
WildcardFilter & DNS name patterns & whether a domain wildcard
  satisfies the specialized filter rule \\
\bottomrule
\end{tabularx}
\end{table*}

The specialization target column names the single construct each template is
written for, which is what keeps it outside the GENERIC set.

\section{Phrase Dictionary of the Mechanical Translation \texorpdfstring{$\smech$}{sigma mech}}
\label{app:sigma}

$\smech$, defined in \S\ref{subsec:mech}, consists of an
atomic-template-to-phrase dictionary $\mathcal{M} : \mathcal{A} \to
\mathcal{L}_{\mathrm{NL}}$ and three combinator reduction rules. Representative
entries appear in Table~\ref{tab:phrasedict}.

\begin{table}[!htbp]
\centering
\footnotesize
\caption{Representative entries of the phrase dictionary $\mathcal{M}$.}
\label{tab:phrasedict}
\begin{tabularx}{\columnwidth}{@{}p{3.1cm} X@{}}
\toprule
Atomic template & $\mathcal{M}(a)$ phrase template \\
\midrule
ExtPresent(O) & ``the \{O\} extension is present'' \\
ExtCritical(O) & ``the \{O\} extension is present and marked
  critical'' \\
FieldEq(F, v) & ``\{F\} equals \{v\}'' \\
FieldEncodedAs(F, T) & ``\{F\} is encoded as ASN.1 \{T\}'' \\
KeyUsageHas(B) & ``the KeyUsage bit \{B\} is asserted'' \\
ListAllMatch(F, T) & ``every entry of \{F\} satisfies
  the result of $\smech(T)$'' \\
SubtreeIPListAny\-HasOctetCount\-AndNotAllZero(F, n) & ``the
  NameConstraints IP subtree \{F\} contains a non-zero entry of \{n\}
  octets'' \\
\bottomrule
\end{tabularx}
\end{table}

\begingroup\emergencystretch=2em\relax
The combinator reduction gives $\smech(\neg t) = $ ``NOT ($\smech(t)$)'';
$\smech(t_1 \wedge t_2) = $ ``($\smech(t_1)$) AND ($\smech(t_2)$)''; $\vee$
likewise; the atomic-template base case $\smech(a) = \mathcal{M}(a)$. For a
conditional pair $(p,q)$, when $p = \neg(p')$, it outputs ``WHEN NOT
($\smech(p')$), THEN $\smech(q)$'', the NEG-PRE template for a negated
precondition; otherwise it outputs ``WHEN ($\smech(p)$), THEN $\smech(q)$''; and
when $p = \varepsilon$ it outputs the single sentence $\smech(q)$.
\par\endgroup


\section{Lint Coverage Alignment Procedure}
\label{app:coverage}

The coverage criterion in \S\ref{subsec:coverage} requires information
beyond a metadata description alone. We therefore first summarize each lint into a reverse IR in Stage-0. For
every lint we pass three inputs to an offline summarizer
$M_{\mathrm{summ}}$ under a constrained prompt: the lint's Go source code, its
Pass/Error test certificates, and the description, citation, and severity
metadata fields. From these we recover the same IR schema our normative rules
use, the core four-tuple $\langle$subject, obligation, predicate, constraint$\rangle$
plus a one-sentence code summary. Three considerations motivate these inputs.
First, the CheckApplies and Execute branches expose the concrete field and
threshold the lint actually tests. Second, the Pass/Error samples fix the
direction of the check, that is, which values it accepts and which it rejects,
together with its applicable scope. Third, casting the lint into our four-tuple
lets each normative rule and lint be compared field by field. Extracting the
full four-tuple, not just the code summary sentence, is what lets Stage-2
compare subject, obligation, predicate, and constraint jointly and reach a
defensible full-versus-partial distinction.

% Algorithm 1 -- zlint coverage-alignment determination for lintable normative rules.
% Included from appendix.tex via % Algorithm 1 -- zlint coverage-alignment determination for lintable normative rules.
% Included from appendix.tex via % Algorithm 1 -- zlint coverage-alignment determination for lintable normative rules.
% Included from appendix.tex via \input{algorithms/alg_coverage}.
%
% Typeset single-column (not figure*) at \scriptsize: the algorithmic lines are
% short, so at full width the right half of most lines was blank and the float was
% billed two columns of page area. Long \Comment texts are placed on their own
% \Statex lines because \Comment is right-aligned and cannot wrap. No step, symbol,
% or comment wording was changed.
\begin{figure}[!htbp]
\raggedright
\setlength{\parskip}{0pt}
\makeatletter
\@ifundefined{c@algorithm}{%
  \newcounter{algorithm}%
  \renewcommand{\thealgorithm}{\arabic{algorithm}}%
}{}%
\makeatother
\refstepcounter{algorithm}
\label{alg:coverage}
\Description{Algorithm for matching each lintable normative rule to candidate zlint checks and aggregating full, partial, and uncovered coverage.}
\hrule width \linewidth height 0.6pt\relax
\nointerlineskip
\noindent\textbf{Algorithm~\thealgorithm.} Algorithmic listing for zlint
coverage alignment determination for individual normative rules.
\par
\nointerlineskip
\hrule width \linewidth height 0.4pt\relax
\nointerlineskip
\scriptsize
\begin{algorithmic}[1]
\Statex \textbf{Input:} lintable normative-rule set $B=\{b_1,\dots,b_m\}$, where each normative rule has
  IR four-tuple $\langle$subject, obligation, predicate, constraint$\rangle$,
  source, and section; zlint lint set $L=\{\ell_1,\dots,\ell_n\}$, where each lint has
  description/citation/severity metadata and Pass/Error tests; offline code
  summarizer $M_{\mathrm{summ}}$, and field-level judge $M_{\mathrm{judge}}$
  with temperature $=0$
\Statex \textbf{Output:} per-normative-rule coverage verdict $\mathrm{verdict}(b)\in\{\text{full, partial,
  none}\}$, covered set $\RLcov$, and uncovered code-generation domain $\RLunc$
\Statex // Offline cached Stage-0: summarize each lint into a reverse
  IR
\ForAll{$\ell \in L$}
  \State $\ell.\mathit{ir} \gets M_{\mathrm{summ}}(\ell.\mathrm{src},
    \ell.\mathrm{tests}, \ell.\mathrm{meta})$
  \Statex \Comment{four-tuple $+$ code summary}
\EndFor
\State pre-split candidate pools by Source: $L_{\mathrm{RFC}}\gets$ RFC lints,\;
  $L_{\mathrm{CABF}}\gets$ CABF lints
\ForAll{$b \in B$}
  \State // Stage-1: deterministic candidate retrieval
  \If{$b.\mathrm{source}=\mathrm{RFC}$}
    \State {$C(b)\gets\{\ell\in L_{\mathrm{RFC}}:$}
    \Statex \hspace{\algorithmicindent}
      {$\textsc{SectionPrefixMatch}(b.\mathrm{section},\ell.\mathrm{citedSections})\}$}
  \ElsIf{$b.\mathrm{source}=\mathrm{CABF}$}
    \State $C(b)\gets L_{\mathrm{CABF}}$ \Comment{CABF sections drift $\Rightarrow$ take all}
  \Else
    \State $C(b)\gets\emptyset$
  \EndIf
  \State // Stage-2: field-level judge, compare four-tuple, best across candidates
  \State $(v,\ell^{*})\gets M_{\mathrm{judge}}(b.\mathit{ir}, C(b))$
  \Statex \Comment{$v\in\{\text{full},\text{partial},\text{none}\}$}
  \State // Stage-3: reject positive matches with conflicting concrete subject families
  \If{$v\in\{\text{full},\text{partial}\} \wedge
      \mathrm{Family}(b.\mathrm{subject})\neq\emptyset \wedge
      \mathrm{Family}(\ell^{*}.\mathrm{subject})\neq\emptyset \wedge
      \mathrm{Family}(b.\mathrm{subject})\neq\mathrm{Family}(\ell^{*}.\mathrm{subject})$}
    \State $v\gets\text{none}$
    \Statex \Comment{different concrete families $\Rightarrow$ wrong-field false positive}
  \EndIf
  \State $\mathrm{verdict}(b)\gets v$
\EndFor
\State // Stage-4: deterministic aggregation, coverage counts full only
\State $\RLcov\gets\{b\in B:\mathrm{verdict}(b)=\text{full}\}$;\quad
  $\RLunc\gets B\setminus \RLcov$
\State \textbf{return} $(\{(b,\mathrm{verdict}(b)):b\in B\},\ \RLcov,\ \RLunc)$
\end{algorithmic}
\hrule width \linewidth height 0.6pt\relax
\end{figure}
.
%
% Typeset single-column (not figure*) at \scriptsize: the algorithmic lines are
% short, so at full width the right half of most lines was blank and the float was
% billed two columns of page area. Long \Comment texts are placed on their own
% \Statex lines because \Comment is right-aligned and cannot wrap. No step, symbol,
% or comment wording was changed.
\begin{figure}[!htbp]
\raggedright
\setlength{\parskip}{0pt}
\makeatletter
\@ifundefined{c@algorithm}{%
  \newcounter{algorithm}%
  \renewcommand{\thealgorithm}{\arabic{algorithm}}%
}{}%
\makeatother
\refstepcounter{algorithm}
\label{alg:coverage}
\Description{Algorithm for matching each lintable normative rule to candidate zlint checks and aggregating full, partial, and uncovered coverage.}
\hrule width \linewidth height 0.6pt\relax
\nointerlineskip
\noindent\textbf{Algorithm~\thealgorithm.} Algorithmic listing for zlint
coverage alignment determination for individual normative rules.
\par
\nointerlineskip
\hrule width \linewidth height 0.4pt\relax
\nointerlineskip
\scriptsize
\begin{algorithmic}[1]
\Statex \textbf{Input:} lintable normative-rule set $B=\{b_1,\dots,b_m\}$, where each normative rule has
  IR four-tuple $\langle$subject, obligation, predicate, constraint$\rangle$,
  source, and section; zlint lint set $L=\{\ell_1,\dots,\ell_n\}$, where each lint has
  description/citation/severity metadata and Pass/Error tests; offline code
  summarizer $M_{\mathrm{summ}}$, and field-level judge $M_{\mathrm{judge}}$
  with temperature $=0$
\Statex \textbf{Output:} per-normative-rule coverage verdict $\mathrm{verdict}(b)\in\{\text{full, partial,
  none}\}$, covered set $\RLcov$, and uncovered code-generation domain $\RLunc$
\Statex // Offline cached Stage-0: summarize each lint into a reverse
  IR
\ForAll{$\ell \in L$}
  \State $\ell.\mathit{ir} \gets M_{\mathrm{summ}}(\ell.\mathrm{src},
    \ell.\mathrm{tests}, \ell.\mathrm{meta})$
  \Statex \Comment{four-tuple $+$ code summary}
\EndFor
\State pre-split candidate pools by Source: $L_{\mathrm{RFC}}\gets$ RFC lints,\;
  $L_{\mathrm{CABF}}\gets$ CABF lints
\ForAll{$b \in B$}
  \State // Stage-1: deterministic candidate retrieval
  \If{$b.\mathrm{source}=\mathrm{RFC}$}
    \State {$C(b)\gets\{\ell\in L_{\mathrm{RFC}}:$}
    \Statex \hspace{\algorithmicindent}
      {$\textsc{SectionPrefixMatch}(b.\mathrm{section},\ell.\mathrm{citedSections})\}$}
  \ElsIf{$b.\mathrm{source}=\mathrm{CABF}$}
    \State $C(b)\gets L_{\mathrm{CABF}}$ \Comment{CABF sections drift $\Rightarrow$ take all}
  \Else
    \State $C(b)\gets\emptyset$
  \EndIf
  \State // Stage-2: field-level judge, compare four-tuple, best across candidates
  \State $(v,\ell^{*})\gets M_{\mathrm{judge}}(b.\mathit{ir}, C(b))$
  \Statex \Comment{$v\in\{\text{full},\text{partial},\text{none}\}$}
  \State // Stage-3: reject positive matches with conflicting concrete subject families
  \If{$v\in\{\text{full},\text{partial}\} \wedge
      \mathrm{Family}(b.\mathrm{subject})\neq\emptyset \wedge
      \mathrm{Family}(\ell^{*}.\mathrm{subject})\neq\emptyset \wedge
      \mathrm{Family}(b.\mathrm{subject})\neq\mathrm{Family}(\ell^{*}.\mathrm{subject})$}
    \State $v\gets\text{none}$
    \Statex \Comment{different concrete families $\Rightarrow$ wrong-field false positive}
  \EndIf
  \State $\mathrm{verdict}(b)\gets v$
\EndFor
\State // Stage-4: deterministic aggregation, coverage counts full only
\State $\RLcov\gets\{b\in B:\mathrm{verdict}(b)=\text{full}\}$;\quad
  $\RLunc\gets B\setminus \RLcov$
\State \textbf{return} $(\{(b,\mathrm{verdict}(b)):b\in B\},\ \RLcov,\ \RLunc)$
\end{algorithmic}
\hrule width \linewidth height 0.6pt\relax
\end{figure}
.
%
% Typeset single-column (not figure*) at \scriptsize: the algorithmic lines are
% short, so at full width the right half of most lines was blank and the float was
% billed two columns of page area. Long \Comment texts are placed on their own
% \Statex lines because \Comment is right-aligned and cannot wrap. No step, symbol,
% or comment wording was changed.
\begin{figure}[!htbp]
\raggedright
\setlength{\parskip}{0pt}
\makeatletter
\@ifundefined{c@algorithm}{%
  \newcounter{algorithm}%
  \renewcommand{\thealgorithm}{\arabic{algorithm}}%
}{}%
\makeatother
\refstepcounter{algorithm}
\label{alg:coverage}
\Description{Algorithm for matching each lintable normative rule to candidate zlint checks and aggregating full, partial, and uncovered coverage.}
\hrule width \linewidth height 0.6pt\relax
\nointerlineskip
\noindent\textbf{Algorithm~\thealgorithm.} Algorithmic listing for zlint
coverage alignment determination for individual normative rules.
\par
\nointerlineskip
\hrule width \linewidth height 0.4pt\relax
\nointerlineskip
\scriptsize
\begin{algorithmic}[1]
\Statex \textbf{Input:} lintable normative-rule set $B=\{b_1,\dots,b_m\}$, where each normative rule has
  IR four-tuple $\langle$subject, obligation, predicate, constraint$\rangle$,
  source, and section; zlint lint set $L=\{\ell_1,\dots,\ell_n\}$, where each lint has
  description/citation/severity metadata and Pass/Error tests; offline code
  summarizer $M_{\mathrm{summ}}$, and field-level judge $M_{\mathrm{judge}}$
  with temperature $=0$
\Statex \textbf{Output:} per-normative-rule coverage verdict $\mathrm{verdict}(b)\in\{\text{full, partial,
  none}\}$, covered set $\RLcov$, and uncovered code-generation domain $\RLunc$
\Statex // Offline cached Stage-0: summarize each lint into a reverse
  IR
\ForAll{$\ell \in L$}
  \State $\ell.\mathit{ir} \gets M_{\mathrm{summ}}(\ell.\mathrm{src},
    \ell.\mathrm{tests}, \ell.\mathrm{meta})$
  \Statex \Comment{four-tuple $+$ code summary}
\EndFor
\State pre-split candidate pools by Source: $L_{\mathrm{RFC}}\gets$ RFC lints,\;
  $L_{\mathrm{CABF}}\gets$ CABF lints
\ForAll{$b \in B$}
  \State // Stage-1: deterministic candidate retrieval
  \If{$b.\mathrm{source}=\mathrm{RFC}$}
    \State {$C(b)\gets\{\ell\in L_{\mathrm{RFC}}:$}
    \Statex \hspace{\algorithmicindent}
      {$\textsc{SectionPrefixMatch}(b.\mathrm{section},\ell.\mathrm{citedSections})\}$}
  \ElsIf{$b.\mathrm{source}=\mathrm{CABF}$}
    \State $C(b)\gets L_{\mathrm{CABF}}$ \Comment{CABF sections drift $\Rightarrow$ take all}
  \Else
    \State $C(b)\gets\emptyset$
  \EndIf
  \State // Stage-2: field-level judge, compare four-tuple, best across candidates
  \State $(v,\ell^{*})\gets M_{\mathrm{judge}}(b.\mathit{ir}, C(b))$
  \Statex \Comment{$v\in\{\text{full},\text{partial},\text{none}\}$}
  \State // Stage-3: reject positive matches with conflicting concrete subject families
  \If{$v\in\{\text{full},\text{partial}\} \wedge
      \mathrm{Family}(b.\mathrm{subject})\neq\emptyset \wedge
      \mathrm{Family}(\ell^{*}.\mathrm{subject})\neq\emptyset \wedge
      \mathrm{Family}(b.\mathrm{subject})\neq\mathrm{Family}(\ell^{*}.\mathrm{subject})$}
    \State $v\gets\text{none}$
    \Statex \Comment{different concrete families $\Rightarrow$ wrong-field false positive}
  \EndIf
  \State $\mathrm{verdict}(b)\gets v$
\EndFor
\State // Stage-4: deterministic aggregation, coverage counts full only
\State $\RLcov\gets\{b\in B:\mathrm{verdict}(b)=\text{full}\}$;\quad
  $\RLunc\gets B\setminus \RLcov$
\State \textbf{return} $(\{(b,\mathrm{verdict}(b)):b\in B\},\ \RLcov,\ \RLunc)$
\end{algorithmic}
\hrule width \linewidth height 0.6pt\relax
\end{figure}


Given these reverse IRs, Algorithm~\ref{alg:coverage} aligns each normative
rule to zlint in four further stages that combine deterministic steps with a
single controlled judge. Stage-1 retrieves candidate lints for a normative rule
by provenance. A normative rule from RFC~5280 is matched only against RFC lints
whose cited section shares a common prefix with the normative rule's section,
a test we write \textsc{SectionPrefixMatch}, since RFC section numbers are
stable. Because CABF section numbers drift across versions, every CABF lint is
taken as a candidate for a CABF normative rule. Stage-2 hands the normative
rule's four-tuple and its candidate set to a field-level judge $M_{\mathrm{judge}}$,
which compares the tuples field by field and returns the best candidate with a
full, partial, or none verdict. Stage-3 is a deterministic false-positive
guard. It maps the subject of the normative rule and of the selected lint to a
field family such as KeyUsage or SubjectAltName via $\mathrm{Family}(\cdot)$,
and downgrades a full or partial verdict to none whenever the two resolve to
different nonempty families, so a KeyUsage normative rule is never credited to
a lint that checks SubjectAltName merely because both mention an extension.
Stage-4 aggregates deterministically. A normative rule is covered if and only
if its verdict is full. The LLM is invoked only
at Stage-0 and Stage-2; candidate
retrieval, the family guard, and aggregation are fully deterministic.


%-------------------------------------------------------------------------------
% Stage-Aware Iterative Verification (SAIV) --- moved out of the main body and
% merged inline here; previously the separate file 07_saiv.tex, now deleted.
%-------------------------------------------------------------------------------
\section{Stage-Aware Iterative Verification (SAIV)}
\label{sec:saiv}

\begingroup\emergencystretch=2em\relax

To coordinate the modules and safeguard their output quality, we propose the SAIV framework.
The basic idea is to predefine two quantifiable optimization
objectives, measure the gap between the current and ideal states by each
objective's computable residual, and when a residual has not reached its
threshold, the residual signal locates the offending module
and the corresponding repair is applied; after re-measuring, the loop proceeds
to the next iteration, until both residuals converge. Alongside the two objectives, two
auxiliary checks audit the lintability determination without contributing to the
objective loss; when a lintability label is suspected to be wrong, SAIV updates
the relevant IR fields and applies the lintability determination again in the next
iteration.

\subsection{Pipeline Structure and Partition}
\label{subsec:saiv-formal}

The whole pipeline chains four modules: extraction, lintability determination,
generation, and verification. Lintability determination is the four-condition
predicate of \S\ref{sec:lintability}, generation is the module restricted to the
DSL code space $\Tv$ of \S\ref{sec:codegen}, and verification is the three-layer
semantic alignment of \S\ref{sec:codegen}. Lintability determination decides a
single normative rule; Algorithm~\ref{alg:saiv} applies its set-level lift, which
strips the noise $\RN$ from $\Rkw$ and sorts the remaining candidates retained as
normative rules into lintable $\RL$ and not-lintable $\RU$, giving the partition
$(\RN, \RL, \RU)$. The loop is driven by the computable signals these modules
produce.

\subsection{Two Iteration Objectives and Two Auxiliary Lintability Checks}
\label{subsec:targets}

SAIV fixes two optimization objectives
for the whole pipeline to approach, and treats the two auxiliary checks
on the lintability determination as repair-routing signals rather than
separate objective terms. Because these checks can change which normative rules enter
$\RL$ and $\RLunc$, they can affect the semantic-equivalence residual in later iterations
even though they add no term to $\Ltot$. The two objectives G1 and G2 each carry
a directly computable residual, $\Lrec$ and $\Lcode$, whose weighted sum is the
total loss $\Ltot$, so convergence is defined by $\Ltot$ alone. S1 and S2 are
checks, not objectives, and add no term to $\Ltot$. They form the two sides of a
single consistency test on the lintability determination, S1 auditing the
not-lintable verdict and S2
the lintable verdict, each firing only to correct a suspected lintability error.
The subsequent subsections revolve around the two objectives, while the two
checks gate the labels those objectives consume.

\noindent\textbf{G1, recall completeness.}
\begin{theorem}[Recall-Completeness Invariant]\label{thm:recall}
If the set-level lintability determination assigns each keyword-recalled candidate in
$\Rkw$ to exactly one of the three classes noise, lintable, and
not-lintable, then recall is conserved from recall to lintability
determination:
\begin{equation}\label{eq:cons}
  \Rkw = \RN + \RL + \RU
\end{equation}
\end{theorem}
\noindent where $\RN$ is noise that contains
RFC~2119 keyword matches rejected as noise by the adjudication procedure,
$\RL$ the lintable normative rules, and $\RU$
the not-lintable normative rules. The lintable class $\RL$ is further split by whether
a normative rule is already implemented by an existing external lint tool,
\begin{equation}\label{eq:split}
  \RL = \RLcov + \RLunc
\end{equation}
where $\RLcov$ are normative rules covered by some lint tool and $\RLunc$ are
not-yet-covered normative rules. The first identity~\eqref{eq:cons} is total conservation
from recall to lintability determination, and the second~\eqref{eq:split} splits the
lintable set so the later objectives can be pursued separately. The uncovered
subset $\RLunc$ is the most important observation domain for generation that
fills the ecosystem gap, while the covered subset $\RLcov$ provides external
evidence of reverse lintability. The residual is the recall-conservation
residual $\Lrec$, a symmetric normalization of the gap in
Theorem~\ref{thm:recall}, computable without manual ground truth, as shown in Tables~\ref{tab:rule-recall} and~\ref{tab:coverage}:
\begin{equation}\label{eq:recall-res}
  \Lrec = 1 - \frac{\min\!\big(\Rkw,\; \RN+\RL+\RU\big)}
  {\max\!\big(\Rkw,\; \RN+\RL+\RU\big)}
\end{equation}

\noindent\textbf{G2, code--normative-rule semantic equivalence for uncovered normative rules.} For the
uncovered subset $\RLunc$, the code summary $\smech(t_r)$, the code summary of
normative rule $r$'s DSL tree $t_r$ computed by $\smech$ of \S\ref{subsec:mech}, is
required to be semantically equivalent to the normative rule text.
The residual $\Lcode$ is the average, over $\RLunc$, of the per-normative-rule correctness
label $\lambda_{\mathrm{code}}(r) = \ind[\mathrm{compile}(\rho(t_r))] \cdot
\mathrm{syn}(\smech(t_r), \normrule(r))$, where $\rho(t_r)$ is the rendered Go
code and $\normrule(r)$ the normative rule text. The structural-completeness score is
$\equiv 1$ in the fixed zlint shell, so it drops out:
\begin{equation}\label{eq:code-res}
  \Lcode = 1 - \frac{1}{|\RLunc|}\sum_{r\in\RLunc}\lambda_{\mathrm{code}}(r)
\end{equation}

\noindent\textbf{S1, reverse lintability of covered normative rules.} The first
auxiliary check audits the not-lintable verdict. If a normative rule has a corresponding
implementation in a lint tool such as zlint, it must be lintable, so the check
counts the violations $\Nviol$ and is cleared when $\Nviol = 0$. $\Nviol$ is
computed by a reverse coverage check that runs the coverage judge over
the not-lintable set $\RU$ and counts the normative rules it rates full or partial; each
such hit is adjudicated as either a lintability false negative and re-extracted to
lintable, or a spurious coverage match and kept not-lintable.

\noindent\textbf{S2, synthesizability of lintable normative rules.} The dual check
audits the lintable verdict. A normative rule judged lintable should be synthesizable
within $\Tv$, so a no-template abstention NT returned for such a
normative rule after generation repair and vocabulary expansion
have both been exhausted is a candidate signal that the lintability determination
mislabeled it.
S2 adds no residual to $\Ltot$; it reuses the generator's existing abstention
marker as a diagnostic. Abstentions that survive both repair channels are held
in a pending set $\NTpend$ and reviewed against the four-condition determination
of \S\ref{sec:lintability}: each is adjudicated either a lintability false
positive, re-derived to not-lintable, or a genuine $\Tv$ expressiveness gap,
moved to the logged set $\NTlog$. S2 clears when $|\NTpend| = 0$. Entries in
$\NTlog$ are reported gaps that do not block termination, though their normative
rules stay in $\RLunc$ and keep contributing to $\Lcode$.

The two objective residuals are weighted into the total loss in
Eq.~\eqref{eq:total}, with $w_R + w_C = 1$:
\begin{equation}\label{eq:total}
  \Ltot = w_R\cdot\Lrec + w_C\cdot\Lcode
\end{equation}
Since recall is conserved, $\Lrec = 0$, $\Ltot$ is governed by $\Lcode$ alone.
The two auxiliary checks stay outside $\Ltot$; their signals are the violation
count $\Nviol$ of S1 and the abstention review of S2 in \S\ref{subsec:algo}.

\begin{table*}[!t]
\centering
\small
\caption{The two SAIV objectives and the two auxiliary lintability checks, with
their signals and repair routing.}
\label{tab:saiv-objectives}
\begin{tabularx}{\textwidth}{@{}L{2.8cm} L{4.2cm} L{4.0cm} >{\raggedright\arraybackslash}X@{}}
\toprule
Item & Signal & Main repair direction & Repair target \\
\midrule
\multicolumn{4}{@{}l}{\emph{Iteration objectives: define $\Ltot$ and
  convergence}}\\
G1 Recall completeness & $\Lrec$, Eq.~\eqref{eq:recall-res} & none &
  accounting invariant only \\
G2 Semantic equivalence & $\Lcode$, Eq.~\eqref{eq:code-res} &
  IR-content, DSL-generation, and verification-judgment repairs &
  IR-content correction / DSL-tree re-synthesis / judgment \& summary
  re-check \\
\midrule
\multicolumn{4}{@{}l}{\emph{Auxiliary lintability checks: audit the
  lintability determination, outside $\Ltot$}}\\
S1 Reverse lintability & $\Nviol$, cleared at $\Nviol=0$, \S\ref{subsec:algo} &
  verification-judgment consistency repair &
  lintability false negatives / spurious coverage matches \\
S2 Synthesizability & $\NTpend$, pending NT abstentions, cleared at
  $|\NTpend|=0$, \S\ref{subsec:algo} &
  diagnose lintability vs. generation gap &
  lintability false positives / missing atoms / $\Tv$ expressiveness gaps
  logged in $\NTlog$ \\
\bottomrule
\end{tabularx}
\end{table*}

Table~\ref{tab:saiv-objectives} gives only the item-to-repair-direction correspondence; the
concrete repair operations are defined in
\S\ref{subsec:repair}.
\textsc{RouteRepair} in Algorithm~\ref{alg:saiv} denotes the routing step that
uses the residuals, lintability-check signals, and failure
traces to select the next applicable repair direction among IR-content,
lintability, generation, and verification repair.

\subsection{Repair Operations}
\label{subsec:repair}

SAIV comprises four kinds of repair operations: IR-content
repair, lintability repair, generation repair, and verification repair. The
next round's operation is selected from the signals above, and a repair is
accepted only when the metric decreases.

\begin{itemize}
\item \textbf{IR-content repair} changes the extraction prompt
  and re-runs IR extraction when the source-level meaning captured by the IR is
  wrong. The LLM may return a corrected IR, denoted REPAIR(ir'), or
  abstain with NO\_FIX.
\item \textbf{Lintability repair} is the same re-extraction
  mechanism restricted to the lintability-gate IR fields: obligation,
  assertion\_subject, enforcement\_phase, and rule\_category. It is invoked when
  S1 indicates that a not-lintable normative rule is actually covered by an existing lint,
  or when S2 indicates that a supposedly lintable normative rule may not be expressible as
  a certificate lint. After re-extraction, SAIV re-runs the lintability
  determination and updates the partition $(\RN, \RL, \RU)$.
\item \textbf{Generation repair} routes by cause. A missing template routes to the
  offline vocabulary-expansion channel, which adds a new atomic
  template to $\mathcal{A}$ after certificate-level execution verification; a wrong IR field, such as too coarse a granularity,
  routes to IR-content repair; and a normative rule that proves not to be lintable
  routes to lintability repair. When the template, IR, and
  lintability label are all correct but the meaning is too complex to express
  within $\Tv$, the normative rule is recorded in the logged gap set $\NTlog$
  rather than repaired as code.
\item \begingroup\sloppy \textbf{Verification repair} enlarges the semantic-equivalence
  judgment's semantic neighborhood, such as negation/affirmation interchange or
  one-to-many decomposition, or corrects the inconsistency between the two
  verdict sources.\par\endgroup
\end{itemize}

\subsection{Iteration Algorithm and Termination Conditions}
\label{subsec:algo}

Algorithm~\ref{alg:saiv} gives the full procedure.

% Algorithm 2 -- Stage-Aware Iterative Verification, abbreviated SAIV.
% Included from the SAIV section of appendix.tex via % Algorithm 2 -- Stage-Aware Iterative Verification, abbreviated SAIV.
% Included from the SAIV section of appendix.tex via % Algorithm 2 -- Stage-Aware Iterative Verification, abbreviated SAIV.
% Included from the SAIV section of appendix.tex via \input{algorithms/alg_saiv}.
%
% Typeset at \small in one column. Long comments and conditions use
% unnumbered continuation lines to preserve the algorithm's steps and wording.
% Start the tall listing in the next column; the explicit break keeps the
% preceding subsection heading with its text instead of stretching the gap.
\newpage
\begin{figure}[H]
\raggedright
\setlength{\parskip}{0pt}
\makeatletter
\@ifundefined{c@algorithm}{%
  \newcounter{algorithm}%
  \renewcommand{\thealgorithm}{\arabic{algorithm}}%
}{}%
\makeatother
\refstepcounter{algorithm}
\label{alg:saiv}
\Description{Algorithm for routing extraction, lintability, generation, and verification repairs until residuals converge or the iteration limit is reached.}
\hrule width \linewidth height 0.6pt\relax
\vspace{1pt}
\nointerlineskip
\noindent\textbf{Algorithm~\thealgorithm.} Algorithmic listing for
Stage-Aware Iterative Verification, abbreviated SAIV.
\par
\vspace{1pt}
\nointerlineskip
\hrule width \linewidth height 0.4pt\relax
\nointerlineskip
\small
\algrenewcommand\algorithmicindent{1em}
\begin{algorithmic}[1]
\Statex \textbf{Input:} Web PKI normative-source corpus $D$, threshold $\theta$, max iterations $K$
\Statex \textbf{Output:} final code set $C^{*}$ and termination status
\State $\Rkw \gets \phi_R(D)$;\quad $(\RN,\RL,\RU)\gets \psi_L(\Rkw)$
\Statex \Comment{$\RU=$ not-lintable}
\State $(\RLcov,\RLunc)\gets \textsc{Coverage}(\RL)$
  \Comment{Algorithm~\ref{alg:coverage}}
\State $C \gets \{\phi_G(r): r\in \RLunc\}$
\State $t\gets 0$
\Repeat
  \State compute objective residuals $\Lrec$~\eqref{eq:recall-res},
    $\Lcode$~\eqref{eq:code-res}; lintability-check signals $\Nviol$ (S1) and the
    pending-abstention set $\NTpend$ (S2)
  \If{$\Ltot<\theta$ \textbf{and} $\Nviol=0$ \textbf{and} $|\NTpend|=0$}
    \State \textbf{break} \Comment{objectives converged, S1 and S2 clear}
  \EndIf
  \State $\mathit{op}\gets \textsc{RouteRepair}(\Lrec,\Lcode,\Nviol,\NTpend,\cdots)$
  \Statex \Comment{repair routing, \S\ref{subsec:repair}}
  \If{$\mathit{op}=\mathcal{P}_R$} \Comment{IR-content repair}
    \ForAll{$r$ routed to $\mathcal{P}_R$}
      \State $\mathit{rep}\gets \mathcal{P}_R(\textsc{FailTrace}(r))$
        \Comment{REPAIR(ir$'$) or NO\_FIX}
      \If{$\mathit{rep}=\textsc{No\_Fix}$ \textbf{or}\newline
        \hspace*{4em}\textsc{ReextractCheck}$(\mathit{rep}.ir',\textsc{text}(r))\neq\emptyset$}
        \State record $r$ in irreducible set $\mathcal{R}_{\mathrm{irred}}$
        \Statex \Comment{kept as residual}
      \ElsIf{\textsc{Retest}$(\mathit{rep}.ir')$ passes}
        \State $\mathrm{IR}(r)\gets \mathit{rep}.ir'$ \Comment{accept}
      \EndIf
    \EndFor
  \Else
    \State apply $\mathit{op}\in\{\mathcal{P}_L,\mathcal{P}_G,\mathcal{P}_V\}$
  \EndIf
  \State update $(\RN,\RL,\RU,\RLunc)$ and $C$
  \State $t\gets t+1$
\Until{$t\ge K$ \textbf{or} no residual decreased, the no-progress condition}
\State \textbf{return} $(C,\ \mathit{status})$, with $\mathit{status}=\textsc{closed}$
  on break, else \textsc{max\_iter} if $t\ge K$ and \textsc{dry} otherwise
\end{algorithmic}
\vspace{1pt}
\hrule width \linewidth height 0.6pt\relax
\end{figure}
.
%
% Typeset at \small in one column. Long comments and conditions use
% unnumbered continuation lines to preserve the algorithm's steps and wording.
% Start the tall listing in the next column; the explicit break keeps the
% preceding subsection heading with its text instead of stretching the gap.
\newpage
\begin{figure}[H]
\raggedright
\setlength{\parskip}{0pt}
\makeatletter
\@ifundefined{c@algorithm}{%
  \newcounter{algorithm}%
  \renewcommand{\thealgorithm}{\arabic{algorithm}}%
}{}%
\makeatother
\refstepcounter{algorithm}
\label{alg:saiv}
\Description{Algorithm for routing extraction, lintability, generation, and verification repairs until residuals converge or the iteration limit is reached.}
\hrule width \linewidth height 0.6pt\relax
\vspace{1pt}
\nointerlineskip
\noindent\textbf{Algorithm~\thealgorithm.} Algorithmic listing for
Stage-Aware Iterative Verification, abbreviated SAIV.
\par
\vspace{1pt}
\nointerlineskip
\hrule width \linewidth height 0.4pt\relax
\nointerlineskip
\small
\algrenewcommand\algorithmicindent{1em}
\begin{algorithmic}[1]
\Statex \textbf{Input:} Web PKI normative-source corpus $D$, threshold $\theta$, max iterations $K$
\Statex \textbf{Output:} final code set $C^{*}$ and termination status
\State $\Rkw \gets \phi_R(D)$;\quad $(\RN,\RL,\RU)\gets \psi_L(\Rkw)$
\Statex \Comment{$\RU=$ not-lintable}
\State $(\RLcov,\RLunc)\gets \textsc{Coverage}(\RL)$
  \Comment{Algorithm~\ref{alg:coverage}}
\State $C \gets \{\phi_G(r): r\in \RLunc\}$
\State $t\gets 0$
\Repeat
  \State compute objective residuals $\Lrec$~\eqref{eq:recall-res},
    $\Lcode$~\eqref{eq:code-res}; lintability-check signals $\Nviol$ (S1) and the
    pending-abstention set $\NTpend$ (S2)
  \If{$\Ltot<\theta$ \textbf{and} $\Nviol=0$ \textbf{and} $|\NTpend|=0$}
    \State \textbf{break} \Comment{objectives converged, S1 and S2 clear}
  \EndIf
  \State $\mathit{op}\gets \textsc{RouteRepair}(\Lrec,\Lcode,\Nviol,\NTpend,\cdots)$
  \Statex \Comment{repair routing, \S\ref{subsec:repair}}
  \If{$\mathit{op}=\mathcal{P}_R$} \Comment{IR-content repair}
    \ForAll{$r$ routed to $\mathcal{P}_R$}
      \State $\mathit{rep}\gets \mathcal{P}_R(\textsc{FailTrace}(r))$
        \Comment{REPAIR(ir$'$) or NO\_FIX}
      \If{$\mathit{rep}=\textsc{No\_Fix}$ \textbf{or}\newline
        \hspace*{4em}\textsc{ReextractCheck}$(\mathit{rep}.ir',\textsc{text}(r))\neq\emptyset$}
        \State record $r$ in irreducible set $\mathcal{R}_{\mathrm{irred}}$
        \Statex \Comment{kept as residual}
      \ElsIf{\textsc{Retest}$(\mathit{rep}.ir')$ passes}
        \State $\mathrm{IR}(r)\gets \mathit{rep}.ir'$ \Comment{accept}
      \EndIf
    \EndFor
  \Else
    \State apply $\mathit{op}\in\{\mathcal{P}_L,\mathcal{P}_G,\mathcal{P}_V\}$
  \EndIf
  \State update $(\RN,\RL,\RU,\RLunc)$ and $C$
  \State $t\gets t+1$
\Until{$t\ge K$ \textbf{or} no residual decreased, the no-progress condition}
\State \textbf{return} $(C,\ \mathit{status})$, with $\mathit{status}=\textsc{closed}$
  on break, else \textsc{max\_iter} if $t\ge K$ and \textsc{dry} otherwise
\end{algorithmic}
\vspace{1pt}
\hrule width \linewidth height 0.6pt\relax
\end{figure}
.
%
% Typeset at \small in one column. Long comments and conditions use
% unnumbered continuation lines to preserve the algorithm's steps and wording.
% Start the tall listing in the next column; the explicit break keeps the
% preceding subsection heading with its text instead of stretching the gap.
\newpage
\begin{figure}[H]
\raggedright
\setlength{\parskip}{0pt}
\makeatletter
\@ifundefined{c@algorithm}{%
  \newcounter{algorithm}%
  \renewcommand{\thealgorithm}{\arabic{algorithm}}%
}{}%
\makeatother
\refstepcounter{algorithm}
\label{alg:saiv}
\Description{Algorithm for routing extraction, lintability, generation, and verification repairs until residuals converge or the iteration limit is reached.}
\hrule width \linewidth height 0.6pt\relax
\vspace{1pt}
\nointerlineskip
\noindent\textbf{Algorithm~\thealgorithm.} Algorithmic listing for
Stage-Aware Iterative Verification, abbreviated SAIV.
\par
\vspace{1pt}
\nointerlineskip
\hrule width \linewidth height 0.4pt\relax
\nointerlineskip
\small
\algrenewcommand\algorithmicindent{1em}
\begin{algorithmic}[1]
\Statex \textbf{Input:} Web PKI normative-source corpus $D$, threshold $\theta$, max iterations $K$
\Statex \textbf{Output:} final code set $C^{*}$ and termination status
\State $\Rkw \gets \phi_R(D)$;\quad $(\RN,\RL,\RU)\gets \psi_L(\Rkw)$
\Statex \Comment{$\RU=$ not-lintable}
\State $(\RLcov,\RLunc)\gets \textsc{Coverage}(\RL)$
  \Comment{Algorithm~\ref{alg:coverage}}
\State $C \gets \{\phi_G(r): r\in \RLunc\}$
\State $t\gets 0$
\Repeat
  \State compute objective residuals $\Lrec$~\eqref{eq:recall-res},
    $\Lcode$~\eqref{eq:code-res}; lintability-check signals $\Nviol$ (S1) and the
    pending-abstention set $\NTpend$ (S2)
  \If{$\Ltot<\theta$ \textbf{and} $\Nviol=0$ \textbf{and} $|\NTpend|=0$}
    \State \textbf{break} \Comment{objectives converged, S1 and S2 clear}
  \EndIf
  \State $\mathit{op}\gets \textsc{RouteRepair}(\Lrec,\Lcode,\Nviol,\NTpend,\cdots)$
  \Statex \Comment{repair routing, \S\ref{subsec:repair}}
  \If{$\mathit{op}=\mathcal{P}_R$} \Comment{IR-content repair}
    \ForAll{$r$ routed to $\mathcal{P}_R$}
      \State $\mathit{rep}\gets \mathcal{P}_R(\textsc{FailTrace}(r))$
        \Comment{REPAIR(ir$'$) or NO\_FIX}
      \If{$\mathit{rep}=\textsc{No\_Fix}$ \textbf{or}\newline
        \hspace*{4em}\textsc{ReextractCheck}$(\mathit{rep}.ir',\textsc{text}(r))\neq\emptyset$}
        \State record $r$ in irreducible set $\mathcal{R}_{\mathrm{irred}}$
        \Statex \Comment{kept as residual}
      \ElsIf{\textsc{Retest}$(\mathit{rep}.ir')$ passes}
        \State $\mathrm{IR}(r)\gets \mathit{rep}.ir'$ \Comment{accept}
      \EndIf
    \EndFor
  \Else
    \State apply $\mathit{op}\in\{\mathcal{P}_L,\mathcal{P}_G,\mathcal{P}_V\}$
  \EndIf
  \State update $(\RN,\RL,\RU,\RLunc)$ and $C$
  \State $t\gets t+1$
\Until{$t\ge K$ \textbf{or} no residual decreased, the no-progress condition}
\State \textbf{return} $(C,\ \mathit{status})$, with $\mathit{status}=\textsc{closed}$
  on break, else \textsc{max\_iter} if $t\ge K$ and \textsc{dry} otherwise
\end{algorithmic}
\vspace{1pt}
\hrule width \linewidth height 0.6pt\relax
\end{figure}


\textbf{Termination conditions.} The loop returns \textsc{closed} when
$\Ltot < \theta$ and both checks clear, $\Nviol = 0$ for S1 and
$|\NTpend| = 0$ for S2; \textsc{max\_iter} at the iteration cap $K$; and
\textsc{dry} when no residual decreases in a round. Adjudicated gaps in $\NTlog$
are reported, not pending, so they do not prevent \textsc{closed}. If every
accepted stage-repair operation is monotone, reducing that stage's error or
leaving it unchanged, then $\Ltot$ is non-increasing each round; the actual
system stops as soon as no residual decreases, together with the hard cap $K$,
so the loop terminates in finitely many rounds.
\par\endgroup
